\section*{Capítulo \ref{ch:b2:solid-liquid-diagrams} --- Diagramas binarios sólido--líquido}

\begin{solution}{exo:b2:solid-liquid-diagrams:1}
A \qty{1300}{\degreeCelsius} el \omterm{def:b2:solid-liquid-diagrams:liquidus}{liquidus} está en 0.541 y el \omterm{def:b2:solid-liquid-diagrams:liquidus}{solidus} en 0.609.
0.40: solo líquido. 0.58: líquido (0.541) y \omterm{def:b2:solid-liquid-diagrams:solid-solution}{disolución sólida} (0.609). 0.70:
solo \omterm{def:b2:solid-liquid-diagrams:solid-solution}{disolución sólida}.
\end{solution}

\begin{solution}{exo:b2:solid-liquid-diagrams:2}
Fracción de sólido $(0.58 - 0.541)/(0.609 - 0.541) = 0.57$: \qty{1.15}{mol} de sólido y
\qty{0.85}{mol} de líquido. Níquel: $1.15 \times 0.609 = \qty{0.70}{mol}$ en el
sólido y $0.85 \times 0.541 = \qty{0.46}{mol}$ en el líquido; total \qty{1.16}{mol}
$= 2.00 \times 0.58$.
\end{solution}

\begin{solution}{exo:b2:solid-liquid-diagrams:3}
A presión fija $v = n + 1 - \varphi$ (sin reacción). (a) $1 + 1 - 2 = 0$.
(b) $2 + 1 - 2 = 1$. (c) Líquido y dos \omterm{def:b2:solid-liquid-diagrams:solid-solution}{disoluciones sólidas}: $2 + 1 - 3 = 0$.
\end{solution}

\begin{solution}{exo:b2:solid-liquid-diagrams:4}
Un descenso regular hasta \qty{232}{\degreeCelsius}, una meseta horizontal mientras el estaño
solidifica y de nuevo un descenso regular, más lento cerca de la temperatura ambiente. En la
meseta el líquido y el sólido coexisten a presión fija, \omterm{def:b2:equilibrium-shifts:variance}{varianza} 0: el calor
retirado lo aporta la cristalización, y la temperatura no puede cambiar
hasta que ha solidificado el último líquido.
\end{solution}

\begin{solution}{exo:b2:solid-liquid-diagrams:5}
Benceno: $\ln x_B = -(9870/8.314)(1/269.6 - 1/278.64) = -0.1429$, $x_B =
0.867$. Naftaleno: $\ln x_N = -(19\,100/8.314)(1/269.6 - 1/353.2) = -2.017$,
$x_N = 0.133$. La suma vale 1.000: \qty{269.6}{K} ($\qty{-3.5}{\degreeCelsius}$) es
la temperatura eutéctica, y el líquido eutéctico tiene 0.133 de naftaleno.
\end{solution}

\begin{solution}{exo:b2:solid-liquid-diagrams:6}
El líquido se enfría hasta el \omterm{def:b2:solid-liquid-diagrams:liquidus}{liquidus} del lado del plomo, donde aparecen cristales de (Pb); el
líquido se enriquece en estaño, hasta el eutéctico a \qty{182.2}{\degreeCelsius}
y 62.1~\%. Allí el líquido restante solidifica en la \omterm{def:b2:solid-liquid-diagrams:eutectic}{mezcla eutéctica}.
Fracción másica de eutéctico: $(30 - 18.91)/(62.13 - 18.91) = 0.26$. Al seguir
enfriando, la solubilidad del estaño en (Pb) disminuye (línea discontinua): un poco de (Sn)
precipita dentro de los cristales de (Pb).
\end{solution}

\begin{solution}{exo:b2:solid-liquid-diagrams:7}
23.3~\% de NaCl en masa: $23.3/76.7 = \qty{0.304}{kg}$ de sal por kilogramo de
agua. A \qty{-25}{\degreeCelsius}, por debajo del eutéctico, no puede existir ningún líquido
sean cuales sean las proporciones: la sal y el hielo permanecen uno junto a otro como sólidos.
\end{solution}

\begin{solution}{exo:b2:solid-liquid-diagrams:8}
$x_{\mathrm B} = 0.5$ está entre E$_1$ y el compuesto: cristaliza primero $\mathrm{AB_2}$,
el líquido se desplaza hacia E$_1$ y allí solidifica en el
eutéctico de A y $\mathrm{AB_2}$; al final, A + $\mathrm{AB_2}$.
$x_{\mathrm B} = 0.9$ está entre E$_2$ y B: cristaliza primero B, y el líquido
alcanza E$_2$; al final, $\mathrm{AB_2}$ + B.
\end{solution}

\begin{solution}{exo:b2:solid-liquid-diagrams:9}
En el borde trasero de la zona el líquido solidifica en un sólido que contiene solo 0.78
veces su fracción molar de cobre: el cobre rechazado se queda en el líquido, que
avanza y se enriquece, de modo que el cobre es barrido hacia el extremo de la barra.
Con una razón cercana a 1, cada pasada solo elimina una fracción de la impureza:
hacen falta muchas pasadas, cada una partiendo de la barra que deja la anterior.
\end{solution}

\begin{solution}{exo:b2:solid-liquid-diagrams:10}
Deducción como en el capítulo. Para un líquido diluido, $\ln x_A \approx -x_B$, $1/T -
1/T^* \approx -\Delta T/T^{*2}$, luego $x_B = \Delta_{\text{fus}}H_A\Delta T/(RT^{*2})$;
con $x_B \approx n_B/n_A = bM_A$, $\Delta T = (RT^{*2}M_A/\Delta_{\text{fus}}H_A)\,b$.
Benceno: $K_f = 8.314 \times 278.64^2 \times 0.07811/9870 = \qty{5.11}{K.kg/mol}$.
\end{solution}

\begin{solution}{exo:b2:solid-liquid-diagrams:11}
Por cada \qty{100}{g}: $19.0/24.31 = \qty{0.782}{mol}$ de Mg y $81.0/207.2 =
\qty{0.391}{mol}$ de Pb; razón $2.00$: $\ce{Mg2Pb}$.
\end{solution}

\begin{solution}{exo:b2:solid-liquid-diagrams:12}
La meseta es proporcional a la masa de líquido eutéctico: la desconocida contiene
$120/300 = 0.40$ de él. Del lado del plomo, $w = 18.91 + 0.40 \times (62.13 -
18.91) = 36.2~\%$ de estaño; del lado del estaño, $w = 96.59 - 0.40 \times (96.59 - 62.13) =
82.8~\%$. El primer cambio de pendiente de la \omterm{def:b2:solid-liquid-diagrams:cooling-curve}{curva de enfriamiento} (más alto para la aleación del 36~\%,
cuyo \omterm{def:b2:solid-liquid-diagrams:liquidus}{liquidus} está en el lado empinado del plomo) o una micrografía (cristales primarios de (Pb) o
de (Sn)) permiten distinguirlas.
\end{solution}

\begin{solution}{pb:b2:solid-liquid-diagrams:1}
\textbf{1.} Plomo \qty{327.5}{\degreeCelsius}, estaño \qty{231.9}{\degreeCelsius}.
\textbf{2.} Por cada \qty{100}{g}: $62.13/118.71 = \qty{0.5234}{mol}$ de estaño y
$37.87/207.2 = \qty{0.1828}{mol}$ de plomo; $x_{\mathrm{Sn}} = 0.5234/0.7062 =
0.741$.
\textbf{3.} Puntos (0, 327.5), (100, 231.9), (62.13, 182.2), (18.91, 182.2),
(96.59, 182.2).
\textbf{4.} Como en la figura: líquido; L + (Pb); L + (Sn); (Pb); (Sn); (Pb) + (Sn).
\textbf{5.} L (62.1~\% Sn) $\longrightarrow$ (Pb) (18.9~\%) + (Sn) (96.6~\%).
\textbf{6.} $v = 2 + 1 - 3 = 0$: la temperatura y las tres composiciones están
fijadas.
\textbf{7.} La \omterm{def:b2:solid-liquid-diagrams:solid-solution}{disolución sólida} rica en plomo (Pb): el plomo disuelve estaño, de modo que el sólido en
equilibrio con el líquido ya contiene estaño, en la proporción que se lee en el
\omterm{def:b2:solid-liquid-diagrams:liquidus}{solidus}.
\textbf{8.} Sigue el \omterm{def:b2:solid-liquid-diagrams:liquidus}{liquidus} hacia el eutéctico, enriqueciéndose en estaño,
hasta el 62.13~\%.
\textbf{9.} Líquido de 62.13~\% y (Pb) de 18.91~\% de estaño.
\textbf{10.} (Pb): $(62.13 - 40)/(62.13 - 18.91) = 0.512$; líquido, 0.488.
\textbf{11.} (Pb) de 18.91~\% y (Sn) de 96.59~\%; (Sn): $(40 - 18.91)/(96.59 -
18.91) = 0.271$; (Pb): 0.729.
\textbf{12.} Constituyente eutéctico, 48.8~\% (el líquido presente justo por encima);
(Pb) primario, 51.2~\%. El (Pb) de la pregunta 11 es la suma de los cristales primarios
y de las láminas de (Pb) del eutéctico.
\textbf{13.} Un cambio de pendiente en el \omterm{def:b2:solid-liquid-diagrams:liquidus}{liquidus}, un descenso más lento y luego una meseta a
\qty{182.2}{\degreeCelsius} que dura 0.488 veces la de la aleación eutéctica.
\textbf{14.} $\ln x_{\mathrm{Sn}} = -(7030/8.314)(1/T - 1/505.1)$.
\textbf{15.} $\ln x_{\mathrm{Bi}} = -(11\,300/8.314)(1/T - 1/544.6)$.
\textbf{16.} $x_{\mathrm{Sn}}(T) + x_{\mathrm{Bi}}(T) = 1$.
\textbf{17.} A \qty{110}{\degreeCelsius}: $0.587 + 0.349 = 0.936$; a
\qty{120}{\degreeCelsius}: $0.621 + 0.382 = 1.003$; a
\qty{130}{\degreeCelsius}: $0.655 + 0.417 = 1.072$. La suma alcanza 1 cerca de
\qty{119.5}{\degreeCelsius}: \qty{120}{\degreeCelsius} con una precisión de un grado.
\textbf{18.} $x_{\mathrm{Bi}} = 0.38$; en masa, $0.38 \times 208.98/(0.38 \times
208.98 + 0.62 \times 118.71) = 0.52$.
\textbf{19.} El modelo da \qty{120}{\degreeCelsius} y 52~\% de bismuto, y el
diagrama evaluado, \qty{138.8}{\degreeCelsius} y 57~\%: unos
\qty{19}{\degreeCelsius} demasiado bajo.
\textbf{20.} El estaño que cristaliza no es puro, sino una disolución que contiene
bismuto: su \omterm{def:b2:chemical-potential:chemical-potential}{potencial químico} es menor que el del estaño puro, el sólido es
más estable y cristaliza a partir del líquido a una temperatura más alta para una
composición del líquido dada. La rama del estaño del \omterm{def:b2:solid-liquid-diagrams:liquidus}{liquidus} sube, y corta la
rama del bismuto más arriba (un líquido no ideal añade su propia corrección).
\textbf{21.} Funde y solidifica a una sola temperatura, la más baja del sistema:
no hay intervalo pastoso durante el cual una unión que se moviera se agrietaría.
\textbf{22.} El plomo es tóxico, y los residuos electrónicos acaban en el medio ambiente.
\textbf{23.} Pueden soldarse componentes sensibles al calor a una temperatura más baja;
pero la unión se reblandece a una temperatura más baja en servicio, y el bismuto la vuelve
frágil.
\textbf{24.} \qty{570}{g} de bismuto y \qty{430}{g} de estaño.
\textbf{25.} El modelo de líquido ideal predice un eutéctico a $\boldsymbol{T_E
\approx \qty{120}{\degreeCelsius}}$, frente a los \qty{138.8}{\degreeCelsius}
evaluados.
\end{solution}
